%=====================================================================
%  ch03 —— 原文 648–862：§2.3 主系列表示
%=====================================================================

\subsection{主系列表示}

\ssszh 本节余下部分回顾 $G$ 的 $\mathbb{F}_q$-点的一个基本表示，即其旗流形上的函数空间。
为此，先作若干准备。

\ssszh  \label{ex:actspace}同前，$H$ 记有限群。若 $H$ 作用在有限集 $X$ 上，就诱导出
$H$ 在 $\Fun(X)$ 上的自然表示。确切地说，给定 $H$ 的元素 $h$ 与 $\Fun(X)$ 中的函数
$\phi$，$h$ 对 $\phi$ 的作用给出元素 $h \cdot \phi$，其定义为
%
\[
(h \cdot \phi)(x) := \phi(h^{-1} x), \quad \text{对 } x \in X.
\]
%
相应的左模结构
%
\begin{equation} \label{e:leftmod}
\Fun(H) \otimes \Fun(X) \rightarrow \Fun(X)
\end{equation}
%
可以这样描述。若把作用记为 $a: H \times X \rightarrow X$，则 \eqref{e:leftmod} 由复合
%
\begin{equation} \label{e:actionfun}
\Fun(H) \otimes \Fun(X) \simeq \Fun(H \times X) \xrightarrow{a_*} \Fun(X).
\end{equation}
%
给出。

\ssszh\label{s:funselfdual} 回忆 $\Fun(X)$ 典范地自对偶。确切地说，记
$\Delta: X \rightarrow X \times X$ 为对角映射、$\pi: X \rightarrow \on{pt}$ 为投影，
则有完美配对
%
\[
\Fun(X) \otimes \Fun(X) \simeq \Fun(X \times X) \xrightarrow{\Delta^!}
\Fun(X) \xrightarrow{\pi_*} \Fun(\on{pt}) \simeq k,
\]
%
具体地说，这不过是「$L^2$」内积
%
\[
f \otimes g \mapsto \underset{x}{\Sigma} \hspace{1mm} f(x) \cdot g(x).
\]
%

\ssszh \label{sss:endos}

特别地，给定有限 $H$-集 $X$ 与 $Y$，有典范等同
%
\[
\begin{gathered}
\on{Hom}_{\on{Rep}(H)^\heartsuit}( \Fun(X), \Fun(Y)) \simeq ( \Fun(X)^* \otimes
\Fun(Y))^{H} \\
\simeq \Fun(X \times Y)^H \simeq \Fun(X \times Y/ H),
\end{gathered}
\]
%
其中上标「$H$」记 $H$-不变子空间，而 $X \times Y / H$ 记 $X \times Y$ 关于 $H$ 的对角
作用的轨道集。

具体地说，给定函数 $K \in \Fun(X \times Y)$，对应的同态 $\Fun(X) \rightarrow \Fun(Y)$
由积分变换
%
\[
   \Fun(X) \xrightarrow{\pi_X^!} \Fun(X \times Y) \xrightarrow{ K \cdot - } \Fun(X \times Y) \xrightarrow{\pi_{Y, *}} \Fun(Y)
\]
%
给出，其中 $\pi_X: X \times Y \rightarrow X$ 与 $\pi_Y: X \times Y \rightarrow Y$ 为投影
映射，$K \cdot -$ 指点乘函数。此外，$K$ 是 $H$-不变的，当且仅当相应的积分变换是
$H$-表示的映射。

\ssszh 我们把上述讨论特殊化到 $G(\mathbb{F}_q)$，其中 $\mathbb{F}_q$ 记有 $q$ 个元素的
有限域。

考虑旗流形 $\sB$，它参数化 $G$ 的 Borel 子群。这是整数环上的射影概形。$G$ 通过共轭
Borel 子群自然地作用在 $\sB$ 上，在 $\mathbb{F}_q$-点上给出作用
%
\[
G(\mathbb{F}_q) \times \sB(\mathbb{F}_q) \rightarrow \sB(\mathbb{F}_q).
\]
%
特别地，如第 \ref{ex:actspace} 小节那样，得到幂幺主系列模
%
\[
\Fun(G(\mathbb{F}_q)) \otimes \Fun(\sB(\mathbb{F}_q)) \rightarrow
\Fun(\sB(\mathbb{F}_q)).
\]
%
此外，它的自同态代数，即 {\em Hecke 代数}，按第 \ref{sss:endos} 小节的方式典范地
等同于
%
\[
\Fun( \sB(\mathbb{F}_q) \times \sB(\mathbb{F}_q) / G(\mathbb{F}_q)).
\]
%
我们要强调：这里考虑的不是相应 $\Fq$-簇上的多项式函数代数，而只是这些簇的有理点所成
有限集上 $k$-值函数的代数。

\ssszh 稍具体些：回忆 $\sB$ 参数化 $G$ 的 Borel 子群。这些子群在 $G$ 中彼此共轭，且自
正规化。因此固定一个 Borel 子群 $B$，就把旗流形等同于齐性空间 $G/B$；在 $\Fq$-点上还有
%
\[
\sB(\mathbb{F}_q) \simeq G(\Fq)/B(\Fq).
\]
%
类似地，$(\sB \times \sB) /G$ 等同于 $B \backslash G / B$，于是有同构
%
\[
\Fun( \sB(\mathbb{F}_q) \times \sB(\mathbb{F}_q) / G(\mathbb{F}_q)) \simeq
\Fun( B(\Fq) \backslash G(\Fq) / B(\Fq)).
\]
%
它把左边的代数结构与右边由卷积给出的代数结构对应起来。

\ssszh 为具体起见，我们较详细地讨论最小的非平凡例子。

\begin{ex} 设 $G = GL(V)$，其中 $V$ 是秩为二的自由阿贝尔群。此时 $\sB$ 是射影直线
$\mathbb{P}(V)$，它参数化 $V$ 的秩一商，并带有 $GL(V)$ 的自然作用。要看这一点，取
$V$ 的一组基 $e_1, e_2$，注意投影 $V \rightarrow V / \langle e_2 \rangle$ 的稳定化子
就是上三角矩阵构成的 Borel 子群
%
\[
B =\begin{pmatrix} * & * \\ & * \end{pmatrix},
\]
%
于是得到所要的与 $G/B$ 的等同。特别地，$\sB(\mathbb{F}_q)$ 有 $q+1$ 个点。

模 $\Fun(\sB(\mathbb{F}_q))$ 分解为两个单模，如下。投影 $\pi: \sB \rightarrow \on{pt}$
给出 $G(\mathbb{F}_q)$-模的（分裂）短正合列
%
\[
0 \rightarrow \on{St} \rightarrow \Fun(\sB(\mathbb{F}_q)) \rightarrow
\Fun(\on{pt}(\mathbb{F}_q)) \simeq k \rightarrow 0,
\]
%
其核 $\on{St}$ 即 Steinberg 模，是单模；理由例如是 Hecke 代数为二维。
\end{ex}

\ssszh 本节内容可概括如下。第一，我们回顾了如何把有限群的群代数构造为卷积代数。
第二，对有限群 $G(\Fq)$，我们回顾了它在旗流形 $\sB(\Fq)$ 上函数中的自然作用。
